\section{Tate compression}\label{sec:compression}

In this section, we apply the Tate compression theorem to $X_R$. 

For $G=C_{p^n}$, set
$\mathcal F_i=\{H:H\le C_{p^i}\}$ and
$\mathcal F_{-1}=\varnothing$. For a genuine $G$-spectrum $X$, put
\[
X^{\Phi_i}=\widetilde E\mathcal F_{i-1}\wedge X,\qquad
X^{h_i}=F((E\mathcal F_i)_+,X^{\Phi_i}),\qquad
X^{t_i}=\widetilde E\mathcal F_i\wedge X^{h_i}.
\]
We use $X^{\Phi_i}$ for $0\le i\le n$ and $X^{h_i},X^{t_i}$ for $0\le i<n$.
Write $\Phi^{C_{p^i}}X$ for geometric fixed points and
$\theta_i:X^{\Phi_{i+1}}\to X^{t_i}$ for the map from isotropy separation~\cite[Section~17]{GreenleesMay}~\cite[Section~2.2]{AMR}.
A spectrum $Y$ is $\mathcal F_i$-complete if
$Y\to F((E\mathcal F_i)_+,Y)$ is an equivalence.

The following form of Tate compression applies to every representation
suspension of $X_R$.

\begin{proposition}\label{prop:compression}
Suppose each $\Phi^{C_{p^i}}X$ has bounded-below underlying spectrum.
For $0\le i<n-1$, $X^{t_i}$ is $\mathcal F_{i+1}$-complete, and
the map $\theta_i$ factors through
$X^{h_{i+1}}$ with contractible space of factorizations.
For $n\ge1$,
\begin{equation}\label{eq:tate-compression}
X\simeq\holim\bigl(
X^{h_0}\to X^{t_0}\leftarrow X^{h_1}
\to\cdots\to X^{t_{n-1}}\leftarrow X^{\Phi_n}\bigr).
\end{equation}
These statements apply to $X=S^V\wedge X_R$ for every $V\in RO(G)$.
\end{proposition}

\begin{proof}
For $r\ge i+1$, the completion map on $C_{p^r}$-fixed points is
\[
(\Phi^{C_{p^i}}X)^{tC_{p^{r-i}}}
\longrightarrow
\bigl((\Phi^{C_{p^i}}X)^{tC_p}\bigr)^{hC_{p^{r-i-1}}},
\]
an equivalence by \cite[Lemma~II.4.1]{NS}.
For $H\in\mathcal F_i$, the $H$-fixed points of the source and target
of the completion map are contractible. Hence
\[
X^{t_i}\xrightarrow{\simeq}
F((E\mathcal F_{i+1})_+,X^{t_i}).
\]
By the universal property of completion,
\[
\Map(X^{h_{i+1}},X^{t_i})\xrightarrow{\simeq}
\Map(X^{\Phi_{i+1}},X^{t_i}).
\]
Equation~\eqref{eq:tate-compression} is
\cite[Theorem~4.7]{ZhangHZ}.
On $G$-fixed points, this gives the iterated pullback of
\cite[Corollary~II.4.7 and Remark~II.4.8]{NS}.
Finally,
\[
\Phi^{C_{p^i}}(S^V\wedge X_R)
\simeq S^{V^{C_{p^i}}}\wedge\Phi^{C_{p^i}}X_R,
\qquad
\pi_k\Phi^{C_{p^i}}(S^V\wedge X_R)=0
\quad(k<\dim_{\mathbb R}V^{C_{p^i}}),
\]
since the geometric fixed points of $X_R$ are connective.
\end{proof}

The vertices of \eqref{eq:tate-compression} are $X^{h_i},X^{\Phi_n}$,
and its edges are $X^{t_i}$. For $n\ge1$, write
$\ell_i:X^{h_i}\to X^{t_i}$ ($0\le i<n$),
$\rho_i:X^{h_{i+1}}\to X^{t_i}$ ($0\le i<n-1$), and
$\rho_{n-1}=\theta_{n-1}:X^{\Phi_n}\to X^{t_{n-1}}$
for the maps in \eqref{eq:tate-compression}. The \emph{incidence map} is
\[
\Delta:\left(\prod_{i<n}X^{h_i}\right)\times X^{\Phi_n}
\longrightarrow\prod_{i<n}X^{t_i},\qquad
(\Delta x)_i=\ell_i(x_i)-\rho_i(x_{i+1}).
\]
Thus $X\simeq\operatorname{fib}(\Delta)$. For $n=0$, the diagram
is $X^{\Phi_0}=X$.

For $V\in RO(G)$, regard $V^{C_{p^i}}$ as a representation
of $C_{p^{n-i}}$ via the $p^i$-power identification. After taking
$G$-fixed points in \eqref{eq:tate-compression} for $S^{-V}\wedge X_R$,
the vertices and edges at $i<n$ are, respectively,
\[
(S^{-V^{C_{p^i}}}\wedge X_R)^{hC_{p^{n-i}}},\qquad
(S^{-V^{C_{p^i}}}\wedge X_R)^{tC_{p^{n-i}}}\quad(0\le i<n),
\]
and the terminal vertex is $S^{-V^{C_{p^n}}}\wedge X_R$. Each $\rho_i$ is induced by
\[
S^{-V^{C_{p^{i+1}}}}\wedge X_R
\simeq\Phi^{C_p}(S^{-V^{C_{p^i}}}\wedge X_R)
\longrightarrow(S^{-V^{C_{p^i}}}\wedge X_R)^{tC_p}.
\]
For $V=0$, this map is cyclotomic Frobenius.
For general $V$, its expression in terms of cyclotomic Frobenius
and the Euler map is \eqref{eq:app-moving-arrow}.

